% Options for packages loaded elsewhere
\PassOptionsToPackage{unicode}{hyperref}
\PassOptionsToPackage{hyphens}{url}
%
\documentclass[
]{article}
\usepackage{amsmath,amssymb}
\usepackage{iftex}
\ifPDFTeX
  \usepackage[T1]{fontenc}
  \usepackage[utf8]{inputenc}
  \usepackage{textcomp} % provide euro and other symbols
\else % if luatex or xetex
  \usepackage{unicode-math} % this also loads fontspec
  \defaultfontfeatures{Scale=MatchLowercase}
  \defaultfontfeatures[\rmfamily]{Ligatures=TeX,Scale=1}
\fi
\usepackage{lmodern}
\ifPDFTeX\else
  % xetex/luatex font selection
\fi
% Use upquote if available, for straight quotes in verbatim environments
\IfFileExists{upquote.sty}{\usepackage{upquote}}{}
\IfFileExists{microtype.sty}{% use microtype if available
  \usepackage[]{microtype}
  \UseMicrotypeSet[protrusion]{basicmath} % disable protrusion for tt fonts
}{}
\makeatletter
\@ifundefined{KOMAClassName}{% if non-KOMA class
  \IfFileExists{parskip.sty}{%
    \usepackage{parskip}
  }{% else
    \setlength{\parindent}{0pt}
    \setlength{\parskip}{6pt plus 2pt minus 1pt}}
}{% if KOMA class
  \KOMAoptions{parskip=half}}
\makeatother
\usepackage{xcolor}
\setlength{\emergencystretch}{3em} % prevent overfull lines
\providecommand{\tightlist}{%
  \setlength{\itemsep}{0pt}\setlength{\parskip}{0pt}}
\setcounter{secnumdepth}{-\maxdimen} % remove section numbering
\usepackage{math-proofread-style}
\EnableMathProofColors
\ifLuaTeX
  \usepackage{selnolig}  % disable illegal ligatures
\fi
\IfFileExists{bookmark.sty}{\usepackage{bookmark}}{\usepackage{hyperref}}
\IfFileExists{xurl.sty}{\usepackage{xurl}}{} % add URL line breaks if available
\urlstyle{same}
\hypersetup{
  hidelinks,
  pdfcreator={LaTeX via pandoc}}

\author{}
\date{}

\begin{document}

\hypertarget{apm-a93j02-ux3c3--1ux207fn1-log-2-via-dominated-convergence}{%
\section{\texorpdfstring{APM a93J02:
\(\sum \mathopen{\mDelimiter{(}}-{\mNumber{1}}\mathclose{\mDelimiter{)}}^{n}\mathbin{\mBridgeOperator{/}}\mathopen{\mDelimiter{(}}n+{\mNumber{1}}\mathclose{\mDelimiter{)}} = \log {\mNumber{2}}\)
via dominated
convergence}{APM a93J02: \textbackslash sum \textbackslash mathopen\{\textbackslash mDelimiter\{(\}\}-\{\textbackslash mNumber\{1\}\}\textbackslash mathclose\{\textbackslash mDelimiter\{)\}\}\^{}\{n\}\textbackslash mathbin\{\textbackslash mBridgeOperator\{/\}\}\textbackslash mathopen\{\textbackslash mDelimiter\{(\}\}n+\{\textbackslash mNumber\{1\}\}\textbackslash mathclose\{\textbackslash mDelimiter\{)\}\} = \textbackslash log \{\textbackslash mNumber\{2\}\} via dominated convergence}}\label{apm-a93j02-ux3c3--1ux207fn1-log-2-via-dominated-convergence}}

\hypertarget{informal-proof}{%
\subsection{1. Informal proof}\label{informal-proof}}

\textbf{Why it's hard.} The geometric series
\({\mNumber{1}}\mathbin{\mBridgeOperator{/}}\mathopen{\mDelimiter{(}}{\mNumber{1}}+x\mathclose{\mDelimiter{)}} = \sum \mathopen{\mDelimiter{(}}-x\mathclose{\mDelimiter{)}}^{n}\)
converges only for \(|x| < {\mNumber{1}}\), yet we need to integrate up
to \(x = {\mNumber{1}}\) (where the series converges conditionally but
not absolutely). Interchanging the sum and integral requires
justification beyond termwise integration of a uniformly convergent
series, since the convergence is not uniform on
\(\mathopen{\mDelimiter{[}}{\mNumber{0}},{\mNumber{1}}\mathclose{\mDelimiter{]}}\).

\textbf{The key insight.} Integrate the \emph{partial sums}
\(S_{n}\mathopen{\mDelimiter{(}}x\mathclose{\mDelimiter{)}} = \sum_{k={\mNumber{0}}}^{n} \mathopen{\mDelimiter{(}}-x\mathclose{\mDelimiter{)}}^{k} = \mathopen{\mDelimiter{(}}{\mNumber{1}}-\mathopen{\mDelimiter{(}}-x\mathclose{\mDelimiter{)}}^{n+{\mNumber{1}}}\mathclose{\mDelimiter{)}}\mathbin{\mBridgeOperator{/}}\mathopen{\mDelimiter{(}}{\mNumber{1}}+x\mathclose{\mDelimiter{)}}\)
over
\(\mathopen{\mDelimiter{[}}{\mNumber{0}},{\mNumber{1}}\mathclose{\mDelimiter{]}}\).
Each partial sum is a polynomial (hence integrable), and
\(\mathopen{\mDelimiter{|}}S_{n}\mathopen{\mDelimiter{(}}x\mathclose{\mDelimiter{)}}\mathclose{\mDelimiter{|}} \le {\mNumber{1}}\mathbin{\mBridgeOperator{/}}\mathopen{\mDelimiter{(}}{\mNumber{1}}-\mathopen{\mDelimiter{|}}\mathopen{\mDelimiter{(}}-x\mathclose{\mDelimiter{)}}\mathclose{\mDelimiter{|}}\mathclose{\mDelimiter{)}}\)
\ldots{} no, we need a uniform dominator. Since
\(\mathopen{\mDelimiter{|}}S_{n}\mathopen{\mDelimiter{(}}x\mathclose{\mDelimiter{)}}\mathclose{\mDelimiter{|}} = \mathopen{\mDelimiter{|}}{\mNumber{1}}-\mathopen{\mDelimiter{(}}-x\mathclose{\mDelimiter{)}}^{n+{\mNumber{1}}}\mathclose{\mDelimiter{|}}\mathbin{\mBridgeOperator{/}}\mathopen{\mDelimiter{|}}{\mNumber{1}}+x\mathclose{\mDelimiter{|}} \le {\mNumber{2}}\mathbin{\mBridgeOperator{/}}\mathopen{\mDelimiter{(}}{\mNumber{1}}+x\mathclose{\mDelimiter{)}} \le {\mNumber{2}}\)
for
\(x \in \mathopen{\mDelimiter{[}}{\mNumber{0}},{\mNumber{1}}\mathclose{\mDelimiter{]}}\),
or more simply: each term
\(\mathopen{\mDelimiter{|}}\mathopen{\mDelimiter{(}}-x\mathclose{\mDelimiter{)}}^{k}\mathclose{\mDelimiter{|}} \le {\mNumber{1}}\),
so
\(\mathopen{\mDelimiter{|}}S_{n}\mathopen{\mDelimiter{(}}x\mathclose{\mDelimiter{)}}\mathclose{\mDelimiter{|}} \le n+{\mNumber{1}}\)
--- not a fixed bound. Better: use
\({\mNumber{1}}\mathbin{\mBridgeOperator{/}}\mathopen{\mDelimiter{(}}{\mNumber{1}}+x\mathclose{\mDelimiter{)}}\)
directly as the dominator. We have
\(S_{n}\mathopen{\mDelimiter{(}}x\mathclose{\mDelimiter{)}} \to {\mNumber{1}}\mathbin{\mBridgeOperator{/}}\mathopen{\mDelimiter{(}}{\mNumber{1}}+x\mathclose{\mDelimiter{)}}\)
pointwise on
\(\mathopen{\mDelimiter{[}}{\mNumber{0}},{\mNumber{1}}\mathclose{\mDelimiter{)}}\),
and
\(\mathopen{\mDelimiter{|}}S_{n}\mathopen{\mDelimiter{(}}x\mathclose{\mDelimiter{)}}\mathclose{\mDelimiter{|}} \le {\mNumber{1}}\mathbin{\mBridgeOperator{/}}\mathopen{\mDelimiter{(}}{\mNumber{1}}-x\mathclose{\mDelimiter{)}}\)
is not integrable. Actually
\(\mathopen{\mDelimiter{|}}\mathopen{\mDelimiter{(}}-x\mathclose{\mDelimiter{)}}^{k}\mathclose{\mDelimiter{|}} = x^{k} \le {\mNumber{1}}\),
so
\(\mathopen{\mDelimiter{|}}S_{n}\mathopen{\mDelimiter{(}}x\mathclose{\mDelimiter{)}}\mathclose{\mDelimiter{|}} \le \mathopen{\mDelimiter{(}}n+{\mNumber{1}}\mathclose{\mDelimiter{)}}\),
which is not uniformly bounded. The right approach is different.

\textbf{Proof.}

\emph{Step 1: Set up the integral.} We have
\(\int_{{\mNumber{0}}}^{{\mNumber{1}}} {\mNumber{1}}\mathbin{\mBridgeOperator{/}}\mathopen{\mDelimiter{(}}{\mNumber{1}}+x\mathclose{\mDelimiter{)}} \mathrm{d}x = \log\mathopen{\mDelimiter{(}}{\mNumber{2}}\mathclose{\mDelimiter{)}}\).

The geometric series gives, for
\(x \in \mathopen{\mDelimiter{[}}{\mNumber{0}},{\mNumber{1}}\mathclose{\mDelimiter{)}}\):
\({\mNumber{1}}\mathbin{\mBridgeOperator{/}}\mathopen{\mDelimiter{(}}{\mNumber{1}}+x\mathclose{\mDelimiter{)}} = \sum_{n={\mNumber{0}}}^{\infty} \mathopen{\mDelimiter{(}}-{\mNumber{1}}\mathclose{\mDelimiter{)}}^{n}x^{n}\).

The partial sums are
\(S_{n}\mathopen{\mDelimiter{(}}x\mathclose{\mDelimiter{)}} = \sum_{k={\mNumber{0}}}^{n} \mathopen{\mDelimiter{(}}-{\mNumber{1}}\mathclose{\mDelimiter{)}}^{k}x^{k}\).
By direct computation:

\(\int_{{\mNumber{0}}}^{{\mNumber{1}}} \mathopen{\mDelimiter{(}}-{\mNumber{1}}\mathclose{\mDelimiter{)}}^{k}x^{k} \mathrm{d}x = \mathopen{\mDelimiter{(}}-{\mNumber{1}}\mathclose{\mDelimiter{)}}^{k}\mathbin{\mBridgeOperator{/}}\mathopen{\mDelimiter{(}}k+{\mNumber{1}}\mathclose{\mDelimiter{)}}\).

So
\(\int_{{\mNumber{0}}}^{{\mNumber{1}}} S_{n}\mathopen{\mDelimiter{(}}x\mathclose{\mDelimiter{)}} \mathrm{d}x = \sum_{k={\mNumber{0}}}^{n} \mathopen{\mDelimiter{(}}-{\mNumber{1}}\mathclose{\mDelimiter{)}}^{k}\mathbin{\mBridgeOperator{/}}\mathopen{\mDelimiter{(}}k+{\mNumber{1}}\mathclose{\mDelimiter{)}}\).

\emph{Step 2: Apply DCT.} We need to show
\(\int_{{\mNumber{0}}}^{{\mNumber{1}}} S_{n}\mathopen{\mDelimiter{(}}x\mathclose{\mDelimiter{)}} \mathrm{d}x \to \int_{{\mNumber{0}}}^{{\mNumber{1}}} {\mNumber{1}}\mathbin{\mBridgeOperator{/}}\mathopen{\mDelimiter{(}}{\mNumber{1}}+x\mathclose{\mDelimiter{)}} \mathrm{d}x.\)

\begin{itemize}
\tightlist
\item
  \textbf{Pointwise convergence:} For each
  \(x \in \mathopen{\mDelimiter{[}}{\mNumber{0}},{\mNumber{1}}\mathclose{\mDelimiter{)}}\),
  \(S_{n}\mathopen{\mDelimiter{(}}x\mathclose{\mDelimiter{)}} \to {\mNumber{1}}\mathbin{\mBridgeOperator{/}}\mathopen{\mDelimiter{(}}{\mNumber{1}}+x\mathclose{\mDelimiter{)}}\)
  (geometric series).
\item
  \textbf{Domination:}
  \(\mathopen{\mDelimiter{|}}S_{n}\mathopen{\mDelimiter{(}}x\mathclose{\mDelimiter{)}}\mathclose{\mDelimiter{|}} = \mathopen{\mDelimiter{|}}{\mNumber{1}}-\mathopen{\mDelimiter{(}}-x\mathclose{\mDelimiter{)}}^{n+{\mNumber{1}}}\mathclose{\mDelimiter{|}}\mathbin{\mBridgeOperator{/}}\mathopen{\mDelimiter{(}}{\mNumber{1}}+x\mathclose{\mDelimiter{)}} \le \mathopen{\mDelimiter{(}}{\mNumber{1}}+x^{n+{\mNumber{1}}}\mathclose{\mDelimiter{)}}\mathbin{\mBridgeOperator{/}}\mathopen{\mDelimiter{(}}{\mNumber{1}}+x\mathclose{\mDelimiter{)}} \le \mathopen{\mDelimiter{(}}{\mNumber{1}}+{\mNumber{1}}\mathclose{\mDelimiter{)}}\mathbin{\mBridgeOperator{/}}\mathopen{\mDelimiter{(}}{\mNumber{1}}+{\mNumber{0}}\mathclose{\mDelimiter{)}} = {\mNumber{2}}\)
  for
  \(x \in \mathopen{\mDelimiter{[}}{\mNumber{0}},{\mNumber{1}}\mathclose{\mDelimiter{]}}\).
  More precisely:
  \(\mathopen{\mDelimiter{|}}S_{n}\mathopen{\mDelimiter{(}}x\mathclose{\mDelimiter{)}}\mathclose{\mDelimiter{|}} = \mathopen{\mDelimiter{|}}\mathopen{\mDelimiter{(}}{\mNumber{1}}-\mathopen{\mDelimiter{(}}-x\mathclose{\mDelimiter{)}}^{n+{\mNumber{1}}}\mathclose{\mDelimiter{)}}\mathbin{\mBridgeOperator{/}}\mathopen{\mDelimiter{(}}{\mNumber{1}}+x\mathclose{\mDelimiter{)}}\mathclose{\mDelimiter{|}} \le {\mNumber{2}}\mathbin{\mBridgeOperator{/}}\mathopen{\mDelimiter{(}}{\mNumber{1}}+x\mathclose{\mDelimiter{)}} \le {\mNumber{2}}\).
  The constant function 2 is integrable on
  \(\mathopen{\mDelimiter{[}}{\mNumber{0}},{\mNumber{1}}\mathclose{\mDelimiter{]}}\).
\end{itemize}

By the dominated convergence theorem:

\(\int_{{\mNumber{0}}}^{{\mNumber{1}}} S_{n}\mathopen{\mDelimiter{(}}x\mathclose{\mDelimiter{)}} \mathrm{d}x \to \int_{{\mNumber{0}}}^{{\mNumber{1}}} {\mNumber{1}}\mathbin{\mBridgeOperator{/}}\mathopen{\mDelimiter{(}}{\mNumber{1}}+x\mathclose{\mDelimiter{)}} \mathrm{d}x.\)

\emph{Step 3: Combine.}

\(\sum_{k={\mNumber{0}}}^{n} \mathopen{\mDelimiter{(}}-{\mNumber{1}}\mathclose{\mDelimiter{)}}^{k}\mathbin{\mBridgeOperator{/}}\mathopen{\mDelimiter{(}}k+{\mNumber{1}}\mathclose{\mDelimiter{)}} = \int_{{\mNumber{0}}}^{{\mNumber{1}}} S_{n}\mathopen{\mDelimiter{(}}x\mathclose{\mDelimiter{)}} \mathrm{d}x \to \int_{{\mNumber{0}}}^{{\mNumber{1}}} {\mNumber{1}}\mathbin{\mBridgeOperator{/}}\mathopen{\mDelimiter{(}}{\mNumber{1}}+x\mathclose{\mDelimiter{)}} \mathrm{d}x = [\log\mathopen{\mDelimiter{(}}{\mNumber{1}}+x\mathclose{\mDelimiter{)}}]_{{\mNumber{0}}}^{{\mNumber{1}}} = \log {\mNumber{2}}.\)

Therefore
\(\sum_{n={\mNumber{0}}}^{\infty} \mathopen{\mDelimiter{(}}-{\mNumber{1}}\mathclose{\mDelimiter{)}}^{n}\mathbin{\mBridgeOperator{/}}\mathopen{\mDelimiter{(}}n+{\mNumber{1}}\mathclose{\mDelimiter{)}} = \log {\mNumber{2}}.\)
∎

\textbf{What connects.} This is a model application of the DCT to
justify termwise integration of a power series at the boundary of its
disk of convergence. The Abel summation theorem gives the same result
(if \(\sum a_{n}\) converges and
\(f\mathopen{\mDelimiter{(}}x\mathclose{\mDelimiter{)}} = \sum a_{n}x^{n}\),
then
\(\sum a_{n} = \lim_{x\to {\mNumber{1}}^{-}} f\mathopen{\mDelimiter{(}}x\mathclose{\mDelimiter{)}}\)),
but the DCT approach is more elementary and directly justifies the
interchange of sum and integral. The dominator
\(\mathopen{\mDelimiter{|}}S_{n}\mathclose{\mDelimiter{|}} \le {\mNumber{2}}\)
comes from the closed-form formula for geometric partial sums, which is
the key algebraic input. The same technique proves
\(\int_{{\mNumber{0}}}^{{\mNumber{1}}} \log\mathopen{\mDelimiter{(}}{\mNumber{1}}+x\mathclose{\mDelimiter{)}}\mathbin{\mBridgeOperator{/}}x \mathrm{d}x = \pi^{{\mNumber{2}}}\mathbin{\mBridgeOperator{/}}{\mNumber{12}}\)
and many other series-integral identities.

\hypertarget{lean-4-theorem-statement}{%
\subsection{2. Lean 4 theorem
statement}\label{lean-4-theorem-statement}}

\begin{verbatim}
import Mathlib.MeasureTheory.Integral.IntervalIntegral
import Mathlib.MeasureTheory.Integral.FundThmCalculus
import Mathlib.Analysis.SpecialFunctions.Log.Basic
import Mathlib.Analysis.SpecificLimits.Normed
import Mathlib.Topology.Algebra.InfiniteSum.Basic

open MeasureTheory Set Filter Topology Real

noncomputable section

/-- Main theorem: Σ (-1)ⁿ/(n+1) = log 2. -/
theorem alternating_harmonic_eq_log_two :
    HasSum (fun n : ℕ => (-1 : ℝ) ^ n / (↑n + 1)) (log 2) := by
  sorry

/-- Equivalent tsum formulation. -/
theorem tsum_alternating_harmonic :
    ∑' n : ℕ, (-1 : ℝ) ^ n / (↑n + 1) = log 2 := by
  sorry

/-- ∫₀¹ 1/(1+x) dx = log 2. -/
lemma integral_inv_one_plus :
    ∫ x in (0:ℝ)..1, 1 / (1 + x) = log 2 := by
  sorry

/-- The partial sum Sₙ(x) = Σₖ₌₀ⁿ (-x)ᵏ = (1-(-x)^{n+1})/(1+x). -/
lemma geom_partial_sum (x : ℝ) (hx : x ≠ -1) (n : ℕ) :
    ∑ k in Finset.range (n + 1), (-x) ^ k =
      (1 - (-x) ^ (n + 1)) / (1 + x) := by
  sorry

/-- ∫₀¹ (-1)ᵏ xᵏ dx = (-1)ᵏ/(k+1). -/
lemma integral_pow_term (k : ℕ) :
    ∫ x in (0:ℝ)..1, (-1 : ℝ) ^ k * x ^ k =
      (-1) ^ k / (↑k + 1) := by
  sorry

/-- ∫₀¹ Sₙ(x) dx = Σₖ₌₀ⁿ (-1)ᵏ/(k+1). -/
lemma integral_partial_sum (n : ℕ) :
    ∫ x in (0:ℝ)..1, ∑ k in Finset.range (n + 1), (-1 : ℝ) ^ k * x ^ k =
      ∑ k in Finset.range (n + 1), (-1) ^ k / (↑k + 1) := by
  sorry

/-- Domination: |Sₙ(x)| ≤ 2 for x ∈ [0,1]. -/
lemma partial_sum_bound (n : ℕ) (x : ℝ) (hx : x ∈ Icc (0:ℝ) 1) :
    |∑ k in Finset.range (n + 1), (-1 : ℝ) ^ k * x ^ k| ≤ 2 := by
  sorry

/-- Pointwise convergence: Sₙ(x) → 1/(1+x) for x ∈ [0,1). -/
lemma partial_sum_tendsto (x : ℝ) (hx : x ∈ Ico (0:ℝ) 1) :
    Tendsto (fun n => ∑ k in Finset.range (n + 1), (-1 : ℝ) ^ k * x ^ k)
      atTop (𝓝 (1 / (1 + x))) := by
  sorry

/-- DCT application: ∫₀¹ Sₙ → ∫₀¹ 1/(1+x). -/
lemma integral_partial_sum_tendsto :
    Tendsto
      (fun n => ∫ x in (0:ℝ)..1,
        ∑ k in Finset.range (n + 1), (-1 : ℝ) ^ k * x ^ k)
      atTop (𝓝 (∫ x in (0:ℝ)..1, 1 / (1 + x))) := by
  sorry

end
\end{verbatim}

\hypertarget{mathlib-cross-references}{%
\subsection{3. Mathlib
cross-references}\label{mathlib-cross-references}}

\textbf{Types/structures:} - \texttt{HasSum\ f\ s} --- the series
\(\sum f\mathopen{\mDelimiter{(}}n\mathclose{\mDelimiter{)}}\) converges
to s - \texttt{tsum} --- the unconditional sum (when HasSum holds) -
\texttt{intervalIntegral} / \(\int x \in a..b, f x\) ---
Riemann/Lebesgue integral on intervals - \texttt{Finset.range\ (n+1)}
--- the finite set \{0, \ldots, n\} for partial sums -
\texttt{Filter.atTop} --- for \(n \to \infty\)

\textbf{Key lemmas:} -
\texttt{MeasureTheory.tendsto\_integral\_of\_dominated\_convergence} ---
DCT: the core tool. Pointwise convergence + integrable dominator ⟹
integral convergence - \texttt{geom\_sum\_eq} /
\texttt{Finset.geom\_sum} ---
\(\sum_{k={\mNumber{0}}}^{n} r^{k} = \mathopen{\mDelimiter{(}}{\mNumber{1}} \mBridgeOperator{-} r^{n+{\mNumber{1}}}\))/(1-r).
In \texttt{Mathlib.Algebra.GeomSum} -
\texttt{hasSum\_geometric\_of\_abs\_lt\_one} ---
\(\sum r^{n} = {\mNumber{1}}\mathbin{\mBridgeOperator{/}}\mathopen{\mDelimiter{(}}{\mNumber{1}}-r\mathclose{\mDelimiter{)}}\)
for \(|r| < {\mNumber{1}}\) -
\texttt{intervalIntegral.integral\_comp\_mul\_deriv} --- for computing
\(\int x^{k} -\) \texttt{integral\_pow} ---
\(\int_{{\mNumber{0}}}^{{\mNumber{1}}} x^{n} \mathrm{d}x = {\mNumber{1}}\mathbin{\mBridgeOperator{/}}\mathopen{\mDelimiter{(}}n+{\mNumber{1}}\mathclose{\mDelimiter{)}}\).
May be in \texttt{Mathlib.Analysis.SpecialFunctions.Integrals} -
\texttt{Real.log\_eq\_integral\_inv} ---
\(\log\mathopen{\mDelimiter{(}}b\mathbin{\mBridgeOperator{/}}a\mathclose{\mDelimiter{)}} = \int_{a}^{b}\)
1/x dx (or related FTC form) - \texttt{intervalIntegral.integral\_sum}
--- interchanging finite sum and integral -
\texttt{HasSum.tendsto\_sum\_nat} --- relating HasSum to the limit of
partial sums - \texttt{abs\_neg\_one\_pow\_mul\_pow\_le} --- bounding
\(\mathopen{\mDelimiter{|}}\mathopen{\mDelimiter{(}}-{\mNumber{1}}\mathclose{\mDelimiter{)}}^{k}x^{k}\mathclose{\mDelimiter{|}}\)
(may not exist by this name)

\textbf{Tactic hints:} - \texttt{apply\ HasSum.of\_tendsto\_sum\_range}
--- converting the limit of partial sums to HasSum -
\texttt{rw\ {[}integral\_partial\_sum,\ integral\_inv\_one\_plus{]}} ---
chain of equalities -
\texttt{apply\ tendsto\_integral\_of\_dominated\_convergence\ (fun\ \_\ =\textgreater{}\ (2:ℝ))}
--- applying DCT with constant dominator 2 -
\texttt{simp\ {[}Finset.sum\_range\_succ{]}} --- unfolding partial sums
- \texttt{norm\_num} --- for computing
\(\mathopen{\mDelimiter{(}}-{\mNumber{1}}\mathclose{\mDelimiter{)}}^{k}\)
and
\({\mNumber{1}}\mathbin{\mBridgeOperator{/}}\mathopen{\mDelimiter{(}}k+{\mNumber{1}}\mathclose{\mDelimiter{)}}\)
- \texttt{ring} --- for algebraic identity
\(\mathopen{\mDelimiter{(}}{\mNumber{1}}-\mathopen{\mDelimiter{(}}-x\mathclose{\mDelimiter{)}}^{n+{\mNumber{1}}}\mathclose{\mDelimiter{)}}\mathbin{\mBridgeOperator{/}}\mathopen{\mDelimiter{(}}{\mNumber{1}}+x\mathclose{\mDelimiter{)}}\)
- \texttt{field\_simp} --- clearing denominators - \texttt{positivity}
--- for \({\mNumber{1}}+x > {\mNumber{0}}\) on
\(\mathopen{\mDelimiter{[}}{\mNumber{0}},{\mNumber{1}}\mathclose{\mDelimiter{]}}\)
- \texttt{gcongr} --- for bounding
\(\mathopen{\mDelimiter{|}}S_{n}\mathopen{\mDelimiter{(}}x\mathclose{\mDelimiter{)}}\mathclose{\mDelimiter{|}} \le {\mNumber{2}}\)
- \texttt{exact?} --- finding the geometric series and integral lemmas

\end{document}
